\documentclass[11pt]{article}
\usepackage[margin=1in]{geometry}
\usepackage{amsmath,amssymb,graphicx,booktabs,hyperref,xcolor}
\title{Computational Replication Report:\\ Jungwirth \emph{et al.} (2025), \emph{Altermagnetic spintronics} (arXiv:2508.09748)}
\author{Ollie (OpenClaw @ Argonne) --- TEXTURES-100 Wave 4}
\date{2026-07-19}
\begin{document}
\maketitle

\section{Paper summary}
This review connects altermagnetism to spintronics. Altermagnets --- $d$-wave (or higher
even-parity) collinear compensated magnets --- enable strongly \emph{spin-polarized electrical
currents without any net magnetization}, plus THz-range spin dynamics, making them attractive for
MRAM-type memory. The conceptual anchor is Fig.~1's model $d$-wave altermagnet, contrasted against
a ferromagnet and a conventional antiferromagnet.

\section{Replication target}
We reproduce Fig.~1's model spintronic signatures via constant-$\tau$ Boltzmann transport on the
$d$-wave spin-split Fermi surfaces:
$\sigma^{s}_{ab}=\sum_{\mathbf k} v^s_a v^s_b\,(-\partial f/\partial E)$ per spin $s$,
charge current $\propto\sigma^{\uparrow}+\sigma^{\downarrow}$, spin current
$\propto\sigma^{\uparrow}-\sigma^{\downarrow}$.

\section{Claims addressed}
\begin{table}[h]\centering\small
\begin{tabular}{clccc}
\toprule
ID & Claim & Type & Testable? & Tested? \\
\midrule
C1 & Current polarization spin-up (bias$\|x$) reverses to spin-down (bias$\|y$) & transport & yes & yes \\
C2 & Diagonal bias: unpolarized longitudinal + transverse pure spin current (SSE) & transport & yes & yes \\
C3 & FM: one polarization all dirs; AFM: unpolarized all dirs & transport & yes & yes \\
C4 & Material TMR/AHE magnitudes; THz dynamics & survey & partly & \textcolor{red}{out of scope} \\
\bottomrule
\end{tabular}
\end{table}

\section{Results vs.\ paper}
\begin{table}[h]\centering\small
\begin{tabular}{lccc}
\toprule
Quantity & Expected (Fig.~1) & This work & Verdict \\
\midrule
$P$ (bias$\|x$) & $+$ (spin-up) & $+0.931$ & \checkmark \\
$P$ (bias$\|y$) & $-$ (spin-down) & $-0.931$ & \checkmark \\
$P_{\rm long}$ (diagonal) & $0$ (unpolarized) & $0.000$ & \checkmark \\
transverse spin current (diag) & $\neq 0$ (SSE) & $\neq 0$ & \checkmark \\
transverse charge current (diag) & $0$ & $0.000$ & \checkmark \\
FM $P$ (x,y,diag) & same sign & $+0.117$ all & \checkmark \\
AFM $P$ (x,diag) & $0$ & $0.000$ & \checkmark \\
\bottomrule
\end{tabular}
\end{table}

\section{Per-claim what worked / what didn't}
\textbf{C1 (worked).} $P(x)=+0.931$, $P(y)=-0.931$: the electrical-current spin polarization
reverses sign with bias direction, exactly the $d$-wave signature. \textbf{C2 (worked).} Diagonal
bias yields zero longitudinal polarization and zero transverse charge current but a nonzero
transverse \emph{spin} current --- the non-relativistic spin-splitter effect. \textbf{C3 (worked).}
The ferromagnet keeps one polarization sign for every bias direction; the conventional AFM
($t_{\rm AM}{=}0$) is unpolarized for all directions --- the two Fig.~1 contrasts. \textbf{C4 (out of
scope).} Absolute TMR/AHE ratios and THz resonance frequencies are review surveys of many papers.

\section{Critique of the replication}
This reproduces the CONCEPTUAL model of a review (Fig.~1), not a single quantitative material
result --- appropriate for a review. Strengths: the transport computation is a genuine
Boltzmann calculation (not hand-waving), and every symmetry signature comes out with the correct
sign and the correct zeros (diagonal, AFM). Honest limitations: (i) the transverse spin current is
reported in raw (unnormalized) simulation units --- only its nonzero-ness and the zero transverse
charge current are the physical statements; (ii) constant-$\tau$ clean limit (no disorder/finite-$T$,
Open Q2); (iii) no junction TMR (Open Q1) or dynamics (Open Q3). LLM-judge: REPLICATED (coverage 9,
agreement 9).

\section{Verdict}
\textbf{REPLICATED} --- the review's model $d$-wave altermagnet spintronic signatures (polarization
reversal, spin-splitter effect, FM/AFM contrasts) reproduced with correct signs and zeros via
Boltzmann transport. Material TMR/AHE magnitudes and THz dynamics out of scope (survey-level).

\section*{Open Questions}
See \texttt{report/open\_questions.json}: Q1 junction TMR magnitude; Q2 disorder/finite-$T$ SSE
robustness; Q3 THz resonance origin; Q4 g-/i-wave transport angular fingerprint; Q5 coupling to
ferroelectricity/superconductivity.

\end{document}
